Compared to the baselines, our results are clearly the ones closest to optimality on both synthetic and real datasets. %The difference between our algorithm and the competitors are significant on synthetic data, and in terms of both $\epsilon_{PEHE}$ on IHDP. %AUUC being an unreliable criterion, we only report those results for the sake of completeness. This metric must be taken with caution because if a model predicts more positive outcomes than existing ones, its AUUC may be higher than the true model.
%We proposed a novel framework for estimating the ITE through the parametric estimation of a mixture of distributions under some causality constraints. In this context, we developed an effective algorithm which has, in addition to the convergence property, the capacity to recover true labels under non informative prior and uniform features assumptions. 
%Injecting the causal constraints in an EM algorithm, we propose an effective and simple algorithm (namely CEM) which has, in addition to the convergence property, the capacity to recover true labels under non informative prior and uniform features assumptions. 
%Compared to other baseline algorithms, our algorithm, using a Gaussian mixture model, performs significantly better on both synthetic and real datasets. 
Moreover, our model is intrinsically more interpretable than the compared baselines as the parameters of the distributions of the causal groups provide information about the causal mechanism at play. 
%We restricted our experiments to a Gaussian mixture but it could be extended to any distribution. %and can even be adapted to Variational Bayesian methods supporting distributions with a marginal likelihood which is difficult to calculate.
%Moreover, our model is flexible and can be adapted to Variational Bayesian methods supporting more distributions~\cite{blei2017variational}, to multiple treatments, to non-compliance to treatments cases~\cite{Imbens1997}, to separate labels~\cite{yamane2018uplift} and to recommendations~\cite{bonner2018causal}.
Finally, our model is versatile and can be adapted to multiple treatments~~\cite{frolich2004programme}, non-compliance to treatment cases~\cite{Imbens1997} or separate labels~\cite{yamane2018uplift}.